% New English translation from checked cote 25, site 54 (PDF page 55).
% AI-translated by GPT-6 Astra (OpenAI, ChatGPT Pro mode).
\scanpage{25/54}
\mspage{\hand{49}}

us, in the present version with L\unsure{,|.}, to have a nec and sufficient condition
\unsure{for|for} geom irreducibility \hand{of $Y_\xi$, $\xi$ gen. in $L$.}

From the cunutesque standpoint, \var{\xout{it}}{}{typewriter strike-out}{ega25_25_054-v0002} and in terms of field
theory, we should express the \var{\xout{(sufficient)}}{}{typewriter strike-out}{ega25_25_054-v0003} condition in terms of
transcendence degree as \unsure{follows|follows}: we choose a “hyperplane
at infinity” containing neither $L^{o}$ nor X, and work in its \unsure{complement|complement},
i.e. essentially on \unsure{the|the} standard affine scheme. We choose
a basis of the \unsure{space|space} of linear forms vanishing on $L^{o}$, \var{\xout{namely}}{}{typewriter strike-out}{ega25_25_054-v0004}
\var{\xout{and consider}}{}{typewriter strike-out}{ega25_25_054-v0005} say $T_{1},\ldots,T_{p}$ (p = codim ($L^{o}$,P)), and consider
their inverse images $t_{1},..,t_{p}$ in the fraction field K of X
(\var{\inserted{X}}{X}{insertion}{ega25_25_054-v0006} assumed integral). At least one of the $t_{i}$, say $t_{1}$, is $\hand{\neq} 0$; let us then
consider $a_{1} = t_{2}/\unsure{t|l}_{1}$,..., $a_{p-1} = t_{p}/t_{1}$. Then dim g(X') is none
other than the transcendence degree of $k(a_{1},..,a_{p-1}) \hand{\subset} K$ over k. Thus
if this transcendence degree is $\hand{\geqslant} 2$, O.K. \var{\xout{Likewise if it}}{}{typewriter strike-out}{ega25_25_054-v0007} If it is 1,
then we must require that over $\overset{\hand{\rule{0.6em}{0.35pt}}}{k}$, f(X) is contained in a linear
subvariety of P containing $L^{o}$ and of dimension one greater\hand{, and that the
generic fiber of $g : X' \to \overline{g(X')}$ is geom. irred.}

Suppose that $L^{o}$ has dimension q; then the fibers of u:
$P - L^{o} \hand{\to} Q$ have dimension q+1, \var{\inserted{hence those of $g$ have dim $\leq q+1$,}}{hence those of $g$ have dim $\leq q+1$,}{insertion}{ega25_25_054-v0008} and consequently we have
\[ \dim g(X') \hand{\geqslant} \dim f(X) - (q+1) , \]
so the dimension condition on g(X') holds provided
\[ \dim f(X) \hand{\geqslant} q+\hand{3} \quad . \]
When q=0, we find the fact mentioned in 8.17 a). \var{\inserted[typed]{\var{\xout{of the}}{}{typewriter strike-out}{ega25_25_054-v0010}}}{\var{\xout{of the}}{}{typewriter strike-out}{ega25_25_054-v0010}}{insertion}{ega25_25_054-v0009} \var{\xout{When q=0,}}{}{typewriter strike-out}{ega25_25_054-v0011}

Returning to the conditions of 8.18, we see that condition a)
concerning an $x \hand{\in} T$ is \unsure{satisfied|satisfied} provided X is not “conical”
at x relative to \unsure{$f$|£}”, in an evident sense.

See whether these last Bertinian developments would fit better
in the next n\textsuperscript{o}, on changing the projective immersion. \var{\margin{\hand{54}}}{\placedmargin{\hand{54}}}{marginal note placed}{ega25_25_054-v0012}
\tnote{25-054-1}{The source's unusual adjective cunutesque is retained without expansion. Its letter doubts in d'/p', façon, complémentaire, le, éspace/espace, and satisfaite remain marked on the corresponding English words. The canceled extra t in suffisante gives no independent English correction. The notation and fiber-dimension claim are retained as written.}

% NOTE: Two small curved marks stand above and below the beginning of canceled Lorsque q=0. A long pencil stroke is in the left margin beside the final paragraph; no deletion is inferred from it.
